\documentclass[11pt]{article}
\IfFileExists{preamble.tex}{% Shared preamble for the Exam 1 study materials
\usepackage[T1]{fontenc}
\usepackage[utf8]{inputenc}
\usepackage{mathpazo}
\linespread{1.08}
\usepackage[margin=1in]{geometry}
\usepackage{amsmath,amssymb,amsthm,bm}
\usepackage{booktabs,array}
\usepackage{enumitem}
\usepackage{xcolor}
\usepackage[most]{tcolorbox}
\usepackage{titlesec}
\usepackage{fancyhdr}
\usepackage[colorlinks=true,linkcolor=navy,urlcolor=navy]{hyperref}

\definecolor{navy}{HTML}{34495E}
\definecolor{slate}{HTML}{5D6D7E}
\definecolor{paleblue}{HTML}{EAF0F6}
\definecolor{palegold}{HTML}{FBF5E6}
\definecolor{gold}{HTML}{B7950B}
\definecolor{palegreen}{HTML}{EAF5EE}
\definecolor{green}{HTML}{1E8449}

\titleformat{\section}{\Large\bfseries\color{navy}}{}{0pt}{}[\vspace{-0.6ex}{\color{navy}\rule{\linewidth}{0.8pt}}]
\titleformat{\subsection}{\large\bfseries\color{navy}}{}{0pt}{}
\titleformat{\subsubsection}{\normalsize\bfseries\color{slate}}{}{0pt}{}
\setcounter{secnumdepth}{0}

\pagestyle{fancy}
\fancyhf{}
\renewcommand{\headrulewidth}{0.4pt}
\fancyhead[L]{\small\color{slate}ECON-GA 1001: Exam 1 Preparation}
\fancyhead[R]{\small\color{slate}\leftmark}
\fancyfoot[C]{\small\thepage}

\setlist[enumerate]{itemsep=3pt,topsep=4pt}
\setlist[itemize]{itemsep=2pt,topsep=3pt}

% Boxes
\newtcolorbox{keybox}[1][Key results]{enhanced,breakable,colback=paleblue,colframe=navy,
  fonttitle=\bfseries,title=#1,boxrule=0.6pt,arc=2pt,left=8pt,right=8pt,top=6pt,bottom=6pt}
\newtcolorbox{planbox}[1][Session plan]{enhanced,breakable,colback=palegold,colframe=gold,
  fonttitle=\bfseries,title=#1,boxrule=0.6pt,arc=2pt,left=8pt,right=8pt,top=6pt,bottom=6pt}
\newtcolorbox{solbox}[1][Answer]{enhanced,breakable,colback=palegreen,colframe=green,
  fonttitle=\bfseries,title=#1,boxrule=0.6pt,arc=2pt,left=8pt,right=8pt,top=6pt,bottom=6pt}

% Problems
\newcounter{prob}[section]
\newcommand{\prob}[2]{\par\medskip\refstepcounter{prob}\noindent{\bfseries\color{navy}Problem \theprob\ \ }{\small\color{slate}[#1]}\ \ \textit{#2}\par\smallskip}
\newcommand{\sol}[1]{\par\medskip\noindent{\bfseries\color{green}Solution #1}\par\smallskip}

% Math shortcuts
\newcommand{\E}{\mathbb{E}}
\DeclareMathOperator{\Var}{Var}
\DeclareMathOperator{\Cov}{Cov}
\DeclareMathOperator{\Corr}{Corr}
\DeclareMathOperator{\SE}{SE}
\DeclareMathOperator{\SD}{SD}
\DeclareMathOperator{\MSE}{MSE}
\DeclareMathOperator{\Bias}{Bias}
\DeclareMathOperator{\rank}{rank}
\DeclareMathOperator{\tr}{tr}
\DeclareMathOperator*{\argmin}{arg\,min}
\DeclareMathOperator{\plim}{plim}
\newcommand{\Col}{\operatorname{Col}}
\newcommand{\bX}{\mathbf{X}}
\newcommand{\by}{\mathbf{y}}
\newcommand{\be}{\bm{\varepsilon}}
\newcommand{\bb}{\bm{\beta}}
\newcommand{\bbh}{\hat{\bm{\beta}}}
\newcommand{\beh}{\hat{\bm{\varepsilon}}}
\newcommand{\iotav}{\bm{\iota}}
\newcommand{\bI}{\mathbf{I}}
\newcommand{\bP}{\mathbf{P}}
\newcommand{\bM}{\mathbf{M}}
\newcommand{\xbar}{\bar{x}}
\newcommand{\ybar}{\bar{y}}
\newcommand{\Xbar}{\bar{X}}
\newcommand{\convp}{\xrightarrow{\;p\;}}
\newcommand{\convd}{\xrightarrow{\;d\;}}
\newcommand{\blank}{\underline{\hspace{3cm}}}
}{% Shared preamble for the Exam 1 study materials
\usepackage[T1]{fontenc}
\usepackage[utf8]{inputenc}
\usepackage{mathpazo}
\linespread{1.08}
\usepackage[margin=1in]{geometry}
\usepackage{amsmath,amssymb,amsthm,bm}
\usepackage{booktabs,array}
\usepackage{enumitem}
\usepackage{xcolor}
\usepackage[most]{tcolorbox}
\usepackage{titlesec}
\usepackage{fancyhdr}
\usepackage[colorlinks=true,linkcolor=navy,urlcolor=navy]{hyperref}

\definecolor{navy}{HTML}{34495E}
\definecolor{slate}{HTML}{5D6D7E}
\definecolor{paleblue}{HTML}{EAF0F6}
\definecolor{palegold}{HTML}{FBF5E6}
\definecolor{gold}{HTML}{B7950B}
\definecolor{palegreen}{HTML}{EAF5EE}
\definecolor{green}{HTML}{1E8449}

\titleformat{\section}{\Large\bfseries\color{navy}}{}{0pt}{}[\vspace{-0.6ex}{\color{navy}\rule{\linewidth}{0.8pt}}]
\titleformat{\subsection}{\large\bfseries\color{navy}}{}{0pt}{}
\titleformat{\subsubsection}{\normalsize\bfseries\color{slate}}{}{0pt}{}
\setcounter{secnumdepth}{0}

\pagestyle{fancy}
\fancyhf{}
\renewcommand{\headrulewidth}{0.4pt}
\fancyhead[L]{\small\color{slate}ECON-GA 1001: Exam 1 Preparation}
\fancyhead[R]{\small\color{slate}\leftmark}
\fancyfoot[C]{\small\thepage}

\setlist[enumerate]{itemsep=3pt,topsep=4pt}
\setlist[itemize]{itemsep=2pt,topsep=3pt}

% Boxes
\newtcolorbox{keybox}[1][Key results]{enhanced,breakable,colback=paleblue,colframe=navy,
  fonttitle=\bfseries,title=#1,boxrule=0.6pt,arc=2pt,left=8pt,right=8pt,top=6pt,bottom=6pt}
\newtcolorbox{planbox}[1][Session plan]{enhanced,breakable,colback=palegold,colframe=gold,
  fonttitle=\bfseries,title=#1,boxrule=0.6pt,arc=2pt,left=8pt,right=8pt,top=6pt,bottom=6pt}
\newtcolorbox{solbox}[1][Answer]{enhanced,breakable,colback=palegreen,colframe=green,
  fonttitle=\bfseries,title=#1,boxrule=0.6pt,arc=2pt,left=8pt,right=8pt,top=6pt,bottom=6pt}

% Problems
\newcounter{prob}[section]
\newcommand{\prob}[2]{\par\medskip\refstepcounter{prob}\noindent{\bfseries\color{navy}Problem \theprob\ \ }{\small\color{slate}[#1]}\ \ \textit{#2}\par\smallskip}
\newcommand{\sol}[1]{\par\medskip\noindent{\bfseries\color{green}Solution #1}\par\smallskip}

% Math shortcuts
\newcommand{\E}{\mathbb{E}}
\DeclareMathOperator{\Var}{Var}
\DeclareMathOperator{\Cov}{Cov}
\DeclareMathOperator{\Corr}{Corr}
\DeclareMathOperator{\SE}{SE}
\DeclareMathOperator{\SD}{SD}
\DeclareMathOperator{\MSE}{MSE}
\DeclareMathOperator{\Bias}{Bias}
\DeclareMathOperator{\rank}{rank}
\DeclareMathOperator{\tr}{tr}
\DeclareMathOperator*{\argmin}{arg\,min}
\DeclareMathOperator{\plim}{plim}
\newcommand{\Col}{\operatorname{Col}}
\newcommand{\bX}{\mathbf{X}}
\newcommand{\by}{\mathbf{y}}
\newcommand{\be}{\bm{\varepsilon}}
\newcommand{\bb}{\bm{\beta}}
\newcommand{\bbh}{\hat{\bm{\beta}}}
\newcommand{\beh}{\hat{\bm{\varepsilon}}}
\newcommand{\iotav}{\bm{\iota}}
\newcommand{\bI}{\mathbf{I}}
\newcommand{\bP}{\mathbf{P}}
\newcommand{\bM}{\mathbf{M}}
\newcommand{\xbar}{\bar{x}}
\newcommand{\ybar}{\bar{y}}
\newcommand{\Xbar}{\bar{X}}
\newcommand{\convp}{\xrightarrow{\;p\;}}
\newcommand{\convd}{\xrightarrow{\;d\;}}
\newcommand{\blank}{\underline{\hspace{3cm}}}
}
\fancyhead[R]{\small\color{slate}Mock Exam}

\begin{document}

\begin{center}
{\LARGE\bfseries\color{navy} Mock Exam 1}\\[6pt]
{\large ECON-GA 1001: Applied Statistics and Econometrics I}\\[4pt]
{\color{slate} Take on Mon Oct 19 \quad$\cdot$\quad Time allowed: 2 hours \quad$\cdot$\quad 100 points}
\end{center}

\begin{planbox}[Instructions]
\begin{itemize}
  \item Closed notes. Non-programmable calculator only. Use the course statistical tables (normal and $t$).
  \item Set a timer for 120 minutes. Suggested pacing: Part A 15 min, Part B 15 min, Part C 85 min, check 5 min.
  \item Show your work on Part C. A correct number with no supporting work earns limited credit.
  \item Grade yourself with \texttt{mock\_exam\_solutions.tex} only after the timer ends.
\end{itemize}
\end{planbox}

%======================================================
\section{Part A: Multiple choice (8 questions, 2 points each)}
\markboth{Part A}{}

\begin{enumerate}[label=\textbf{A\arabic*.},leftmargin=*]
\item Which is \textbf{not} a requirement for $f(x)$ to be a valid probability density function?
  \begin{enumerate}[label=(\alph*)]
  \item $f(x)\ge0$ for all $x$ \quad (b) $\int_{-\infty}^{\infty}f(x)\,dx=1$ \quad (c) $f(x)\le1$ for all $x$ \quad (d) $P(a\le X\le b)=\int_a^bf(x)\,dx$
  \end{enumerate}

\item If $\Cov(X,Y)=0$, then
  \begin{enumerate}[label=(\alph*)]
  \item $X$ and $Y$ are independent \quad (b) $\E[Y\mid X]$ is constant \quad (c) $X$ and $Y$ have no linear association \quad (d) $\Var(X+Y)=\Var(X)\Var(Y)$
  \end{enumerate}

\item An estimator has $\E[\hat\theta_n]=\theta$ and $\Var(\hat\theta_n)=\sigma^2/4$ for every $n$. It is
  \begin{enumerate}[label=(\alph*)]
  \item unbiased and consistent \quad (b) biased and consistent \quad (c) unbiased and not consistent \quad (d) biased and not consistent
  \end{enumerate}

\item A 95\% confidence interval for $\mu$ is $[4.2,\ 7.8]$. At the 5\% level, a two-sided test of $H_0:\mu=8$
  \begin{enumerate}[label=(\alph*)]
  \item rejects \quad (b) fails to reject \quad (c) cannot be decided without $n$ \quad (d) accepts $H_0$
  \end{enumerate}

\item In OLS with an intercept, which is \textbf{not} guaranteed by the algebra of least squares?
  \begin{enumerate}[label=(\alph*)]
  \item $\sum_i\hat e_i=0$ \quad (b) $\bX'\beh=\mathbf0$ \quad (c) $\E[\be\mid\bX]=\mathbf0$ \quad (d) $\bar{\hat y}=\ybar$
  \end{enumerate}

\item Which model \textbf{cannot} be written as linear in the parameters?
  \begin{enumerate}[label=(\alph*)]
  \item $Y=\beta_0+\beta_1\ln X+\varepsilon$ \quad (b) $\ln Y=\beta_0+\beta_1X+\varepsilon$ \quad (c) $Y=\beta_0+e^{\beta_1X}+\varepsilon$ \quad (d) $Y=\beta_0+\beta_1X+\beta_2X^2+\varepsilon$
  \end{enumerate}

\item Let $\bP=\bX(\bX'\bX)^{-1}\bX'$ and $\bM=\bI-\bP$. Then $\bP\bM\by$ equals
  \begin{enumerate}[label=(\alph*)]
  \item $\hat\by$ \quad (b) $\beh$ \quad (c) $\mathbf0$ \quad (d) $\by$
  \end{enumerate}

\item $Y$ is regressed on $X$ with an intercept. $Y$ is then measured in cents instead of dollars ($Y^*=100Y$). Which statement is true?
  \begin{enumerate}[label=(\alph*)]
  \item The slope and $R^2$ both rise by a factor of 100.
  \item The slope rises by a factor of 100 and $R^2$ is unchanged.
  \item The slope is unchanged and the residuals rise by a factor of 100.
  \item Nothing changes.
  \end{enumerate}
\end{enumerate}

%======================================================
\section{Part B: Brief answers (4 questions, 4 points each)}
\markboth{Part B}{}

Briefly define or describe. Use mathematical expressions where appropriate.

\begin{enumerate}[label=\textbf{B\arabic*.},leftmargin=*]
\item Estimator versus estimate.
\item Full column rank of $\bX$, and why OLS needs it.
\item Homoskedasticity versus heteroskedasticity. Write the conditional variance-covariance matrix of $\be$ in each case (assume no correlation across observations).
\item The difference between $\E[\be\mid\bX]=\mathbf0$ and $\bX'\beh=\mathbf0$.
\end{enumerate}

%======================================================
\section{Part C: Problems (68 points)}
\markboth{Part C}{}

\subsection{C1. A continuous density (14 points)}
Let $f(x)=c(2x+1)$ for $0\le x\le2$, and $0$ otherwise.
\begin{enumerate}[label=(\alph*)]
\item (3) Find $c$.
\item (3) Derive the CDF $F(x)$ for all real $x$.
\item (2) Find $P(X>1)$.
\item (4) Find $\E[X]$ and $\Var(X)$.
\item (2) Find the median.
\end{enumerate}

\subsection{C2. Joint and conditional distributions (14 points)}
\begin{center}
\begin{tabular}{c|ccc}
$X\backslash Y$ & 0 & 1 & 2 \\ \hline
1 & 0.10 & 0.20 & 0.10 \\
2 & 0.05 & 0.15 & 0.40
\end{tabular}
\end{center}
\begin{enumerate}[label=(\alph*)]
\item (3) Find the marginal distributions of $X$ and $Y$, and $\E[Y]$.
\item (4) Find $\E[Y\mid X=1]$ and $\E[Y\mid X=2]$, and verify the Law of Iterated Expectations.
\item (4) Find $\Cov(X,Y)$.
\item (3) Are $X$ and $Y$ independent? Justify with probabilities.
\end{enumerate}

\subsection{C3. Inference for a mean (12 points)}
A random sample of $n=16$ from a normal population has $\Xbar=52.4$ and $s=4.8$.
\begin{enumerate}[label=(\alph*)]
\item (4) Test $H_0:\mu=50$ against $H_1:\mu\ne50$ at the 5\% level. State the statistic, its reference distribution, the critical value, and your conclusion in words.
\item (3) Construct the 95\% confidence interval and relate it to part (a).
\item (2) Test $H_0:\mu=50$ against $H_1:\mu>50$ at the 5\% level.
\item (3) Suppose instead the population standard deviation is known to be $\sigma=4.8$. Redo part (a) and explain why the reference distribution changes.
\end{enumerate}

\subsection{C4. Simple OLS and rescaling (14 points)}
\begin{center}
\begin{tabular}{c|ccccc}
$x_i$ & 1 & 2 & 4 & 5 & 8 \\ \hline
$y_i$ & 3 & 5 & 6 & 9 & 12
\end{tabular}
\end{center}
\begin{enumerate}[label=(\alph*)]
\item (5) Compute $\hat\beta_1$ and $\hat\beta_0$ for $y_i=\beta_0+\beta_1x_i+e_i$.
\item (3) Compute the residuals and verify that $\sum\hat e_i=0$ and $\sum x_i\hat e_i=0$.
\item (3) Compute TSS, RSS, and $R^2$.
\item (3) Let $y^*_i=10y_i$ and $x^*_i=x_i-2$. Give the new slope, intercept, and $R^2$ without re-running the regression. Explain.
\end{enumerate}

\subsection{C5. Matrix OLS and fit (14 points)}
Consider $\by=\bX\bb+\be$ with $\bX$ an $n\times K$ matrix whose first column is $\iotav$.
\begin{enumerate}[label=(\alph*)]
\item (5) Set up the least squares problem and derive the normal equations. State the assumption needed to solve them, and give $\bbh$.
\item (3) Show that the normal equations imply $\bX'\beh=\mathbf0$ and that $\sum_i\hat e_i=0$.
\item (3) Partition $\bX=[\bX_1\ \ \bX_2]$. Describe, in three steps, how the Frisch-Waugh-Lovell theorem recovers $\bbh_2$. Write $\bM_{X_1}$ explicitly.
\item (3) A regression with $n=40$ and $K=5$ has $R^2=0.60$. Compute $\bar R^2$. Explain why researchers report $\bar R^2$ alongside $R^2$.
\end{enumerate}

\vfill
\begin{center}\color{slate}\itshape End of exam. Stop the timer.\end{center}

\end{document}
